% Round 18. Author: general_u2 branch. Independent review is required.
\section{Numerical local exactness with arbitrary ambient Picard rank}
\label{r18pg:section}

\paragraph{Scope.}
The ambient variety in this section need not be a complete intersection,
and its numerical divisor group need not have rank one.  The conclusion
is unconditional numerical $U_2$-vanishing for the specified product
smoothings.  The Green conclusion additionally uses the stated
curve-comparison hypothesis.  In dimension four an ambient variety
with algebraic geometric $H^2$ supplies that hypothesis and gives full
local G.  No middle Hodge-standard conjecture on a fourfold component
and no general numerical-to-homological comparison are used.

\subsection{The native geometric datum}

Let $R$ be the valuation ring of a finite extension $K/\mathbf Q_p$,
with its original absolute value, uniformizer $\pi$, and finite
residue field $k$.  Let $\mathcal Y/R$ be smooth and projective,
with geometrically integral fibers of dimension $n+1$, where $n\ge4$.
Let $\mathcal L$ be a relatively very ample line bundle and put
$Y=\mathcal Y_k$ and $h=c_1(\mathcal L_k)$.
Choose $m\ge3$, positive integers $d_1,\ldots,d_m$, and native sections
\[
 f_i\in H^0(\mathcal Y,\mathcal L^{d_i}),\qquad
 Q\in H^0(\mathcal Y,\mathcal L^{D}),\qquad D=\sum_i d_i.
 \tag{PG18.1}\label{r18pg:sections}
\]
Write $H_i=V(\bar f_i)$ and $H_I=\bigcap_{i\in I}H_i$.
Require the following actual geometric conditions.
\begin{enumerate}
\item The $H_i$ form a geometric simple-normal-crossings divisor on
      $Y$.  Every $H_I$ with $|I|\le n+1$ is smooth of the expected
      dimension $n+1-|I|$.  Positive-dimensional $H_I$ are geometrically
      integral; zero-dimensional $H_I$ are finite etale.  Intersections
      of $n+2$ factors are empty.
\item $\bar Q$ cuts every positive-dimensional $H_I$ transversely and
      avoids every zero-dimensional $H_I$.  Each positive-dimensional
      $H_I\cap V(\bar Q)$ is geometrically integral.  A zero-dimensional
      such intersection is retained with all of its native points.
\item The generic hypersurface
      \[
        X=V(f_1\cdots f_m+\pi Q)\subset\mathcal Y_K
        \tag{PG18.2}\label{r18pg:generic}
      \]
      is smooth and geometrically integral.
\end{enumerate}
The positive powers of an ample class make every expected
positive-dimensional factor intersection nonempty.  In particular
every pair and every triple of distinct indices occurs.  Native
connected pieces are always retained; there is no division by a
geometric component count.

Perform the ordered ambient blowups along the strict transforms of
the special divisors $H_1,\ldots,H_{m-1}$, and take the strict
transform $\mathcal X$ of \eqref{r18pg:generic}.  The HP18/NF18
chart calculation applies inside the smooth ambient $\mathcal Y$:
at a point on $s$ factors where $\bar Q=0$ the actual etale equation
is $x_1\cdots x_s=\pi y$; each blowup uses the next ideal
$(t,x_i)$ in the preceding chart.  The saturated charts are
\[
 \pi=x_1\cdots x_{a-1}t_{a-1},\qquad
 x_a=t_{a-1}u_a,\qquad
 y=u_a x_{a+1}\cdots x_s.
 \tag{PG18.3}\label{r18pg:charts}
\]
At points where $\bar Q$ is a unit the family was already a strict
crossing and these operations restrict to isomorphisms.  Thus
$\mathcal X$ is regular, integral, flat and projective, with reduced
strictly semistable special fiber $\bigcup_iD_i$, and no extra
vertical component.  The native closed stratum $D_I$ is the iterated
blowup of $H_I$ in the reverse-order centers
\[
 Z_{I,j}=H_I\cap H_j\cap V(\bar Q),\qquad
                j=m,m-1,\ldots,\max I+1.
 \tag{PG18.4}\label{r18pg:centers}
\]
Every occurring center is regular of codimension two and its strict
transform is isomorphic to its original center.  Indeed preceding
centers induce on it a Cartier ideal; at zero-dimensional centers
those preceding intersections are empty by the avoidance hypothesis.
These are the actual glued ideals $(\bar f_j,\bar Q)$ on $H_I$.
They do not depend on the transition units of the line bundles.
All these maps lie over $Y$, so $h$ denotes a single global pullback
class on the resolved model and on its strata.  It is generally nef,
not ample, on a blown-up stratum; the argument below retains this
distinction.

\subsection{The rational Picard comparison used on the starting strata}

\begin{lemma}[Rational Picard Lefschetz in the needed characteristic]
\label{r18pg:picard}
Let $V$ be smooth projective over a field of characteristic $p>0$,
of dimension at least four, and let $E\subset V$ be a smooth ample
Cartier divisor.  Restriction induces an isomorphism
\[
                  \operatorname{Pic}(V)_{\mathbf Q}
                    \xrightarrow{\sim}\operatorname{Pic}(E)_{\mathbf Q}.
 \tag{PG18.5}\label{r18pg:picard-Q}
\]
In the datum above, restriction therefore identifies
\[
 M:=N^1(Y)_{\mathbf Q}\xrightarrow{\sim}N^1(H_i)_{\mathbf Q}
                  \xrightarrow{\sim}N^1(H_i\cap H_j)_{\mathbf Q}.
 \tag{PG18.6}\label{r18pg:ambient-NS}
\]
These are native rational maps, and the assertions reapply over
every finite extension of $k$.
\end{lemma}

\begin{proof}
The rational assertion in positive characteristic can be obtained
directly from SGA~2, Expose XII, Corollaire~3.6, without invoking
Kodaira vanishing.  Since $\dim E\ge3$, Serre duality and Serre
vanishing give an integer $N$ such that
\[
 H^a(E,\mathcal O_E(-bE))=0\qquad(a=1,2,\ b\ge N).
\]
Let $E_N$ be the Cartier thickening $NE$.  For each $s\ge1$, the
sheaf $\mathcal O_{E_N}(-sNE)$ has a filtration with quotients
$\mathcal O_E(-(sN+j)E)$ for $0\le j<N$.  Its first and second
cohomology therefore vanish.  The depth and parafactorial hypotheses
of the cited Corollaire hold on the smooth $V$.  It consequently gives
$\operatorname{Pic}(V)\simeq\operatorname{Pic}(E_N)$.

For successive thickenings, the kernel of the map of unit sheaves
is the additive sheaf $\mathcal O_E(-jE)$.  The cohomology sequence
shows that the kernel and cokernel of
$\operatorname{Pic}(E_{j+1})\to\operatorname{Pic}(E_j)$ are killed
by $p$.  For the finite composite from $E_N$ to $E$ they are killed
by a power of $p$.  Tensoring with $\mathbf Q$ proves
\eqref{r18pg:picard-Q}.  The argument takes place over the original
field; it does not replace native line bundles by geometric ones.

Apply this twice, first from $Y$ to $H_i$, then from $H_i$ to
$H_i\cap H_j$.  The intermediate dimensions are at least four.
It gives surjectivity on the numerical divisor quotients.  For
injectivity, if a native divisor $v$ on the larger variety restricts
numerically to zero, pair its restriction with the restrictions of
all native divisors and the appropriate power of $h$.  Since the
cut divisor has class $d_i h$ or $d_jh$, projection formula gives
a nonzero positive multiple of the divisor Hodge-index form on the
larger variety.  That form is nondegenerate, so $v=0$ numerically.
This proves both numerical isomorphisms.  The same proof, including
the $p$-power thickening step, works after every finite constant
extension.
\end{proof}

Put
\[
 q_Y(\alpha,\beta)=\deg_Y(\alpha\beta h^{n-1}),\qquad
 M_0=\{\alpha\in M:\deg_Y(\alpha h^n)=0\}.
 \tag{PG18.7}\label{r18pg:ambient-form}
\]
Divisor Hodge index gives the rational orthogonal decomposition
\[
                         M=\mathbf Qh\mathbin{\perp}M_0,
              \qquad q_Y|_{M_0}\text{ is negative definite over }\mathbf R.
 \tag{PG18.8}\label{r18pg:ambient-primitive}
\]
No assertion is made about $N^2(Y)$ or $N^2(H_i)$.

\subsection{Actual triple columns and a nef numerical Lefschetz map}

Use the DQ17 signs: $\rho$ is terminal restriction minus initial
restriction and $\sigma$ is initial Gysin minus terminal Gysin.
On an edge $E_{ij}$ with $i<j$, let $q_{ij}:E_{ij}\to H_i\cap H_j$
be the specified composition of blowups and let $e_{ij,l}$ be the
total-transform exceptional divisor for $l>j$.  The Chow blowup
formula, followed by numerical coefficient extraction, gives
\[
       N^1(E_{ij})_{\mathbf Q}
          =q_{ij}^*(M|_{H_i\cap H_j})
                   \oplus\bigoplus_{l>j}\mathbf Qe_{ij,l}.
 \tag{PG18.9}\label{r18pg:edge-basis}
\]
The center of each edge blowup has dimension $n-3\ge1$ and is
geometrically integral by hypothesis, so it contributes exactly
one exceptional divisor.  The maps descending to these numerical
groups are proper pushforward, refined pullback and multiplication
by the stated Chern classes; their preservation of numerical
equivalence follows from the native degree projection formula.

For $i<j<l$, the alternating Gysin of the fundamental class of the
triple $D_{ijl}$ has its three nonzero divisor entries
\[
            \tau_{ijl}=(d_lh_{ij}-e_{ij,l},\ -d_jh_{il},\ d_ih_{jl}).
 \tag{PG18.10}\label{r18pg:triple-columns}
\]
This is an actual Chow identity.  On the first edge its factor
divisor has class $d_lh$ and its own blowup subtracts $e_{ij,l}$.
The two opposite charts at that step separate this strict divisor
from the later centers; hence no further exceptional term appears.
On the other two edges it is the full factor-divisor pullback.
This is precisely the two-parameter chart computation of NF18,
which only uses the factor and $Q$ as local parameters and applies
unchanged in $\mathcal Y$.  In particular the exceptional-coordinate
block of all these columns is $-I$.

Literal opposite pushes of a triple to each component cancel in
Chow, giving $\sigma_2\tau_{ijl}=0$.  Thus every $w\in W^1$ can
be changed, without changing $\sigma_2w$, to a tuple
\[
                        w_{ij}=q_{ij}^*(\alpha_{ij}|_{H_i\cap H_j}),
                         \qquad\alpha_{ij}\in M.
 \tag{PG18.11}\label{r18pg:base-tuple}
\]
This elimination does not identify $M$ with its ample line.

\begin{lemma}[The ambient nef class remains injective on edge divisors]
\label{r18pg:nef-injective}
Multiplication by the actual global class $h$ is injective
\[
                             h:W^1\longrightarrow W^2.
 \tag{PG18.12}\label{r18pg:nef-map}
\]
\end{lemma}

\begin{proof}
On an edge use the form
$q_{E,h}(u,v)=\deg_E(uvh^{n-3})$ on numerical divisors.
The starting block under \eqref{r18pg:ambient-NS} is
$d_id_jq_Y$, hence is nondegenerate.  At one codimension-two
blowup in $Z$, projection gives zero cross terms between its
new exceptional divisor and pullback divisors.  If
$j:F\hookrightarrow\widetilde E$, $p:F\to Z$ and
$\xi=c_1(\mathcal O_F(1))$, self-intersection $j^*j_*=-\xi$
and $p_*\xi=1$ give the exceptional square
\[
                             -\deg_Z(h^{n-3})<0.
 \tag{PG18.13}\label{r18pg:exceptional-square}
\]
The last inequality holds because the center is isomorphic to its
original positive-dimensional ample complete intersection and $h$
restricts to its ample class.  Iterate, using total transforms for
the earlier divisors: the new exceptional line is orthogonal to
the entire earlier pullback block.  The final form is thus the
orthogonal sum of $d_id_jq_Y$ and the nonzero negative entries
\eqref{r18pg:exceptional-square}; it is nondegenerate.

If $hu=0$ in $N^2(E)$, then for every divisor $v$ the number
$\deg_E(uvh^{n-3})$ vanishes, by multiplying the assumed zero
class by $vh^{n-4}$.  The exponent is nonnegative because $n\ge4$.
Nondegeneracy forces $u=0$, proving the injection on every edge.
The proof has not called the nef class ample and has not applied
the ample Hodge-index theorem to a blowup pullback class.
\end{proof}

\subsection{The ambient primitive energy removes the non-scalar directions}

\begin{theorem}[Arbitrary-Picard-rank numerical local exactness]
\label{r18pg:U2}
For every model satisfying \eqref{r18pg:sections}--\eqref{r18pg:generic}
and the stated transversality hypotheses,
\[
       U_2=
       \frac{\ker(\rho_2:C^2\to W^2)\cap
                         \operatorname{im}(\sigma_2:W^1\to C^2)}
                     {\operatorname{im}A_1}=0.
 \tag{PG18.14}\label{r18pg:U2zero}
\]
This holds in every $n\ge4$ and imposes no bound on $\dim M$.
True triple and higher strata are retained.
\end{theorem}

\begin{proof}
Let $z=\sigma_2w\in\ker\rho_2$ and put $d=\rho_2\sigma_2$.
Use \eqref{r18pg:triple-columns} to eliminate all exceptional
divisor coordinates of $w$.  The new representative is
\eqref{r18pg:base-tuple} and still has $dw=0$.
Decompose uniquely
\[
             \alpha_{ij}=a_{ij}h+\alpha^0_{ij},
                               \qquad\alpha^0_{ij}\in M_0.
 \tag{PG18.15}\label{r18pg:alpha-decomposition}
\]
Let $w^0$ be the corresponding base-pullback tuple of the
$\alpha^0_{ij}$.  DG18's actual triple factorization
$d=\tau_2\delta_1$, paired with $w^0h^{n-4}$, gives
\[
 \begin{aligned}
 0&=\sum_e\deg_{E_e}(w^0_e\,(dw)_e h^{n-4})\\
  &=\sum_{i<j<l}\deg_{D_{ijl}}
                \bigl((\delta_1w^0)_{ijl}(\delta_1w)_{ijl}h^{n-4}\bigr)\\
  &=\sum_{i<j<l}d_id_jd_l\,
              q_Y(\alpha^0_{ij}-\alpha^0_{il}+\alpha^0_{jl},
                   \alpha^0_{ij}-\alpha^0_{il}+\alpha^0_{jl}).
 \end{aligned}
 \tag{PG18.16}\label{r18pg:primitive-energy}
\]
Here the restrictions of each base-pullback entry to a resolved
triple are the pullbacks of the same ambient divisor.  Its
intersection number is unchanged by the specified proper birational
blowups of that triple.  On the original triple the fundamental
class in $Y$ is $d_id_jd_lh^3$.  These two facts prove the last
equality, and orthogonality $q_Y(M_0,h)=0$ removes exactly the
scalar terms.  No sign assertion for the full surface divisor
space on a triple has been used.

Every weight in the last sum is positive and $q_Y$ is negative
definite on $M_0$.  Consequently
\[
                      \alpha^0_{ij}-\alpha^0_{il}+\alpha^0_{jl}=0
                                \quad(i<j<l).
 \tag{PG18.17}\label{r18pg:cocycle}
\]
Because every pair and triple occurs, this is the ordinary
complete-simplex cocycle equation with coefficient space $M_0$.
It is solved explicitly over $\mathbf Q$ by
$v_1=0$, $v_i=\alpha^0_{1i}$ for $i>1$:
$\alpha^0_{ij}=v_j-v_i$.  Choose native divisor representatives
of the $v_i$ on $Y$ and pull them to $D_i$.  Their actual
restrictions to every edge are the corresponding ambient
pullbacks.  Subtracting their $\rho_1$ image therefore changes
$z$ by an element of $\operatorname{im}A_1$ and leaves a
representative of the form $hb$, with $b\in W^0$.

It remains to remove this scalar part.  Its compatibility says
\[
                   0=\rho_2\sigma_2(hb)=h\rho_1\sigma_1b.
 \tag{PG18.18}\label{r18pg:scalar-step}
\]
Lemma~\ref{r18pg:nef-injective} implies
$\rho_1\sigma_1b=0$.  The unconditional $U_1=0$ of DG18 gives
$\sigma_1b=A_0c$ for some rational $c\in C^0$.  That theorem
uses any relatively ample class on the model and does not require
the present nef $h$ to be its polarization.  Finally the actual
global projection formula gives
$\sigma_2(hb)=hA_0c=A_1(hc)$.
Thus the original $z$ belongs to $\operatorname{im}A_1$, proving
\eqref{r18pg:U2zero}.
\end{proof}

\begin{proposition}[The larger sufficient quotient is genuinely nonzero]
\label{r18pg:Q-dimension}
Write $b=\dim_{\mathbf Q}N^1(Y)_{\mathbf Q}$ and
$c=\binom{m-1}{2}$.  Then
\[
                      \operatorname{rank}(\rho_2\sigma_2)=bc.
 \tag{PG18.23}\label{r18pg:d-rank}
\]
For every relatively ample global rational class $H$ on the
resolved model, the quotient of DG18 satisfies
\[
 Q_H=\frac{W^1}{\operatorname{im}\rho_1+HW^0+\ker\sigma_2},
              \qquad \dim Q_H=(b-1)c.
 \tag{PG18.24}\label{r18pg:Q-rank}
\]
In particular, if $b>1$, that stronger sufficient
divisor-generation criterion fails for every such $H$, even
though the entire $U_2$ vanishes.
\end{proposition}

\begin{proof}
The exceptional block $-I$ in \eqref{r18pg:triple-columns}
identifies $W^1/\operatorname{im}\tau$ with one copy of $M$
per edge.  The map $d=\rho_2\sigma_2$ kills the triple image.
Its kernel on these base representatives is precisely the
space of $M$-valued vertex coboundaries.
Indeed the proof of Theorem~\ref{r18pg:U2} first forces the
$M_0$ part of a kernel representative to be a coboundary.
After subtracting it, $d(hb)=h d_0b=0$ with
$d_0=\rho_1\sigma_1$.  The injection
\eqref{r18pg:nef-map} forces $d_0b=0$.
The scalar triple identity gives
\[
 \sum_e b_e\deg_{E_e}((d_0b)_eh^{n-2})
           =\sum_T\deg_T(h^{n-2})(\delta_0b)_T^2.
\]
All triple degrees here are positive: $h$ pulls back an ample
class from the original triple and its top degree is unchanged
by the specified blowups.  Thus $\ker d_0=\ker\delta_0$.
Conversely $\delta_0b=0$ makes $d_0b=0$ by the same actual
triple factorization.  In the complete two-skeleton,
$\ker\delta_0=\operatorname{im}\rho_0$, of dimension $m-1$.
Hence the remaining scalar tuple is also a vertex coboundary.
All $M$-valued coboundaries are in $\ker d$ by
$\rho_2A_1=0$.  It follows that
\[
 \operatorname{rank}d
       =b\left(\binom m2-(m-1)\right)=bc.
\]

For any global relatively ample $H$, the projection formula
gives $d(Hb)=H d_0b$.  Divisor Hodge index on every edge
makes $H:W^1\to W^2$ injective for $n\ge4$, exactly as in
DG18.  Thus $\operatorname{rank}d(HW^0)=\operatorname{rank}d_0=c$.
The DG18 exact sequence
$0\to U_2\to Q_H\to\operatorname{im}d/d(HW^0)\to0$,
together with $U_2=0$, now gives
$\dim Q_H=bc-c$, proving the assertion.
\end{proof}

This calculation distinguishes the new step from the scalar
generation proof.  The kernel equation $dw=0$ forces the
non-scalar ambient part to satisfy the triple cocycle equation.
Those triple equations remove the nonzero quotient
\eqref{r18pg:Q-rank}; declaring that quotient zero would miss
the mechanism.

\subsection{The separate representative gate and the actual Green conclusion}

\begin{corollary}[Codimension two, and the complete fourfold group]
\label{r18pg:green}
Assume additionally that every numerically zero rational curve cycle
on each native double $E_{ij}$ has zero geometric $\ell$-adic cycle
class, for one $\ell\ne p$.  Then every specified geometrically
homologically trivial rational codimension-two cycle on $X$ satisfies
strong G.  In dimension $n=4$ this gives full local G in every
codimension.

A sufficient additional condition, valid uniformly after finite
constant extensions, is that
\[
     H^2(Y_{\bar k},\mathbf Q_\ell(1))
          \text{ is spanned by geometric divisor classes}.
 \tag{PG18.19}\label{r18pg:H2-algebraic}
\]
Under \eqref{r18pg:H2-algebraic}, the Green assertions apply afresh
to all cycles over every unramified finite valued extension of $K$.
More generally the same extension assertion holds if the stated
curve comparison is supplied on the doubles over every finite
constant extension.  Comparison supplied only over the original
native $k$ gives the Green assertion only over the original $K$.
\end{corollary}

\begin{proof}
First verify the sufficient comparison.  Etale weak Lefschetz for
the two smooth ample cuts identifies $H^2$ of the starting edge
$H_i\cap H_j$ with $H^2(Y)$.  The successive codimension-two
blowup formula adds the geometric $H^0$ of each center, represented
by its exceptional divisors.  Under \eqref{r18pg:H2-algebraic},
the entire geometric $H^2(E_{ij})(1)$ is therefore spanned by
divisors defined over some finite extension of $k$.

A native numerically zero curve stays numerically zero after that
finite extension: its degree against a divisor there is, up to the
positive full field degree, its original native degree against the
properly pushed divisor.  All these degrees are zero.  Poincare
duality between the curve-class degree and $H^2$ now makes its
geometric cycle class zero.  This proves the precise comparison,
including native nonsplit divisor classes; it does not assert a
cycle-class map out of an arbitrary numerical higher cycle group.

CL17 gives $K_{n-2}=B_{n-2}$ from that comparison.  Theorem
\ref{r18pg:U2} gives $U_2=0$.  The DQ17 exact sequence makes
$V_{n-2}/(K_{n-2}+L_{n-2})$ zero.  ER17 and VT17 consequently
kill the obstruction of the original specified cycle.  The frozen
reconstruction produces an ordinary actual delta-form on the
complement of its precise support and proves the Green identity
against every original compact ordinary actual test.  The finite
model tests are not substituted for this final quantifier.

For $n=4$, the other degrees are the frozen divisor, one-cycle and
zero-cycle cases.  This gives the full cycle domain on that fourfold.
An unramified extension preserves the supplied equations, geometric
transversality and the uniformizer exponent.  The Picard and numerical
arguments are repeated over its native residue field.  Under
\eqref{r18pg:H2-algebraic} the geometric divisor-spanning hypothesis
also remains true, so the cohomological comparison is proved anew
there.  The same is valid when uniform comparison over all finite
constant extensions is a separate hypothesis.  New cycles are then
covered directly and are not assumed to descend as pullbacks.
The original native curve-comparison hypothesis alone has not been
shown to ascend to new cycles over a residue extension.  If the
separate RI18 ramification theorem is activated, its
unramified-to-all-finite implication applies to the full fourfold
conclusion under either of the explicitly uniform hypotheses just
stated.  That implication is an additional dependency, not part of
the chart calculation above.
\end{proof}

\subsection{A Picard-rank-two fourfold outside the projective-space CI class}

\begin{corollary}[A concrete deep-stratum example]
\label{r18pg:product-example}
Take $n=4$ and
\[
      \mathcal Y=\mathbf P^1_R\times\mathbf P^4_R,
              \qquad\mathcal L=\mathcal O(1,1),
              \qquad m=5,\qquad d_i=2.
 \tag{PG18.20}\label{r18pg:product-data}
\]
Whenever the stated transversality and generic smoothness conditions
hold, the resulting smooth bidegree-$(10,10)$ fourfold has full local
G, has $U_2=0$ while $\dim Q_H=6$ for every global relatively ample
$H$, and has
\[
                    \operatorname{rank}N^4=159
               \quad\text{on }H^4(X_{\bar K},\mathbf Q_\ell).
 \tag{PG18.21}\label{r18pg:monodromy-example}
\]
Such native data exist at suitable finite places.  Its geometric
numerical Picard rank is two, so it is not isomorphic over any
field extension to any smooth four-dimensional complete intersection
in projective space.  It has no potential good reduction and cannot
have a strict semistable smooth projective dominant witness of
crossing depth less than four.
\end{corollary}

\begin{proof}
The two factor hyperplane classes span the geometric $H^2$ of $Y$,
so Corollary~\ref{r18pg:green} supplies full local G.
Here $b=2$ and $c=\binom42=6$, so
Proposition~\ref{r18pg:Q-dimension} gives $\dim Q_H=6$.
For existence,
choose five sufficiently general rational bidegree-$(2,2)$ sections
on $\mathbf P^1_{\mathbf Q}\times\mathbf P^4_{\mathbf Q}$.
The finite transversality conditions are nonempty open conditions
by characteristic-zero Bertini for these very ample systems.
Choose a rational bidegree-$(10,10)$ section $Q$ cutting every
positive-dimensional stratum smoothly, avoiding all deepest points,
and itself smooth on the ambient variety.  These choices are
possible simultaneously in a nonempty rational parameter open.
Positive-dimensional ample complete intersections are geometrically
connected; together with smoothness they are geometrically integral.

Spread all data and smoothness conditions over $\mathbf Z[1/S]$.
The pencil $\prod f_i+tQ$ has a smooth member at infinity, so its
singular parameters form a proper closed finite subset of the
parameter line in characteristic zero.  For sufficiently large
primes $p\notin S$ with the rational parameter $t=p$ outside that
finite set, take $K=\mathbf Q_p$ and $\pi=p$.  All the native special
and generic hypotheses above hold.

Write $a,b$ for the two factor hyperplane classes.  Then
$\deg_Y(a+b)^5=5$, since $a^2=0$, $b^5=0$, and
$\deg_Y(ab^4)=1$.  Each fivefold factor intersection therefore
has $5\cdot2^5=160$ geometric points.  Every lower factor stratum
is geometrically connected, and the prescribed blowups preserve
this.  Thus the weight-zero restriction complex has the complete
simplex incidence through degree three and has 160 copies of its
unique top four-simplex.  Its last differential has rank one.
The HM18/NG18 extreme weight calculation gives an isomorphism for
$N^4$ between the opposite extreme pieces, of dimension $160-1$.
The same weight bounds and $E_2$ degeneration used there show that
this is the actual rank of $N^4$, not merely a lower bound.  No
general weight-monodromy conjecture is used.  NM17 gives the stated
non-potential-good-reduction and crossing-depth consequences, with
depth defined as the maximum simultaneous component count minus one.

For the geometric distinction, weak Lefschetz for the smooth ample
fourfold $X\subset\mathbf P^1\times\mathbf P^4$ gives $b_2(X)=2$.
The restrictions of $a,b$ are numerically independent: their divisor
intersection matrix with $(a+b)^2$ on $X$ is
\[
                          10\begin{pmatrix}0&1\\1&3\end{pmatrix},
 \tag{PG18.22}\label{r18pg:Picard-matrix}
\]
whose determinant is nonzero.  Divisor numerical/homological
comparison therefore gives geometric numerical Picard rank exactly
two.  A smooth complete-intersection fourfold in projective space has
geometric $b_2=1$, generated by its nonzero hyperplane class, and hence
has numerical Picard rank one.  An isomorphism preserves that rank.
This excludes every such complete intersection, including every smooth
projective-space hypersurface, after arbitrary field extension.

These distinctions do not classify all possible motivic overlaps with
earlier sectors.  The claim is the explicit new supplied model and
cycle-domain certificate, with the stated intrinsic exclusions.
\end{proof}

\paragraph{References, dependencies and remaining questions.}
The primary Picard input is SGA~2, Expose XII, Corollaire~3.6,
printed p.~153 in the original pagination, reproduced on PDF p.~129
of \url{https://arxiv.org/pdf/math/0511279}.  Its cohomology hypotheses
are verified above on the thickening; no characteristic-zero vanishing
theorem has been used in characteristic $p$.
The weak-Lefschetz and Poincare-duality inputs are the standard etale
theorems already used with their precise degree ranges in NF18/CA18.
The smooth complete-intersection rank-one comparison used only for
the final exclusion is recorded by Ghorpade--Lachaud,
\url{https://arxiv.org/abs/0808.2169}, Propositions~3.1--3.2.
The geometric blowup calculation is HP18/NF18 in an arbitrary smooth
ambient chart; the numerical argument uses DQ17, DG18, and ordinary
divisor Hodge index.  All actual Green conclusions pass through
CL17/ER17/VT17 and the unchanged frozen reconstruction.

The general model need not admit this common ambient divisor frame
or this complete-simplex arrangement.  Even here, without the specified
curve-comparison gate, the result is $U_2=0$ and does not by itself
annihilate the CL17 representative defect.  In dimensions greater than
four the present argument proves codimension-two G under that gate,
not full all-degree G.  No global height statement follows without
checking the other finite places.  General G and the unrestricted
complete Beilinson--Bloch height domain remain open.  No priority claim
is made for the classical inputs or for this internally developed
application.
